\chapter{Troubleshooting}
\label{app:trouble}

\section{A method is missing from the binding}

Read \code{bindings/coverage.json} first. It was written before anything was
compiled, and it names every member that did not reach a backend, with a reason.
The likely entries:

\begin{center}
\small
\begin{tabular}{@{}lL{9.0cm}@{}}
\toprule
\textbf{Reason} & \textbf{What to do} \\
\midrule
\code{overload\_not\_expressible} &
  The target keys methods by name and this overload lost the tie. Wrap it in an
  \field{extensions} entry under a distinct name, or reach it through the
  dynamic model. \S\ref{sec:extensions} \\
\addlinespace
\code{unmarshalable\_signature} &
  \code{detail} names the offending type. If it is a foreign container or an
  Eigen type, Chapter~\ref{ch:types} is the fix. If it is a
  \code{std::function} or a raw pointer, write an \field{extensions} wrapper
  with a marshallable signature. \\
\addlinespace
\code{no\_identifier} &
  An operator or conversion function. There is no name to bind to --- expose it
  as a named extension method. \\
\addlinespace
\code{function\_template} &
  A member template has no fixed parameter pack to splice. Add an explicit
  non-template wrapper. \\
\addlinespace
\code{hidden\_by\_derived} &
  C++ name hiding, honoured deliberately. Add \code{using Base::f;} to the
  derived class if you want the base overloads visible. \\
\addlinespace
\code{sequence\_not\_adaptable} &
  A registered sequence with no \code{std::vector} boundary adapter on this
  backend. \\
\addlinespace
\code{extension\_has\_no\_member\_pointer} &
  This backend dispatches through member pointers and cannot take extensions.
  Expected on \lang{qt}, \lang{qml}, \lang{csharp}, \lang{java}. \\
\bottomrule
\end{tabular}
\end{center}

If the member is not in \code{coverage.json} at all, it was not public, or its
class is not in \field{classes}.

\section{The generated driver does not compile}

\subsection*{\code{no member named 'vec3' in namespace 'geom'}}

A glob entry derived a class name from a file stem that names no such class.
See the table in Section~\ref{sec:globs}: list the class explicitly, or add the
header to \field{exclude}.

\subsection*{An error inside a \field{signature} you wrote}

The signature is spliced verbatim. Spell the types the way the header does,
fully qualified. \code{rosetta\_gen} only checks its shape --- the types are the
compiler's business.

\subsection*{An error inside a \field{module\_init} statement}

Same cause: those are spliced verbatim into the module entry point. Spell types
qualified, and remember there are no arguments available at that point.

\subsection*{The bound header sees a different class than the binding will}

A configuration macro is missing from \field{compile\_definitions}. Those are
emitted for the \emph{generator driver} as well as the binding targets precisely
so the two cannot disagree --- if you passed the macro some other way, move it
into the manifest.

\section{A side-car annotation had no effect}

\begin{itemize}
  \item Did you \textbf{re-run \code{rosetta\_gen}}? The JSON is baked into
    \code{bindings.h} at generation time. A rebuild of the binding alone will
    not pick it up.
  \item A key naming no member \emph{fails the build}, so a silent no-effect
    usually means the annotation reached a backend that ignores it ---
    \code{label} and \code{widget} are UI-only. See
    Appendix~\ref{app:annotations}.
  \item \code{widget::slider} without a \code{range} falls back to the default
    editor. \code{widget::radio} without a \code{combobox} does nothing.
\end{itemize}

\section{The binding is slow}

\begin{itemize}
  \item Check the build type. Omitting \field{build\_type} now defaults to
    \code{Release}, but a \code{-DCMAKE\_BUILD\_TYPE=Debug} on the command line
    wins. An unoptimised pybind11 call costs roughly eight times an optimised
    one.
  \item If you are going through the dynamic model, every access is a linear
    name lookup plus a scoring pass. Cache the resolved \code{MetaMethod}, and
    keep an expanded binding for hot paths. \S\ref{ch:dynamic}
  \item A foreign container crosses \emph{by copy} through a
    \code{std::vector} boundary. For bulk data on the Python family, prefer
    \field{interop} --- the numpy caster avoids the copy where the layout
    allows.
\end{itemize}

\section{The module will not load}

\begin{itemize}
  \item \textbf{Missing shared library.} Every \field{user\_lib} \field{dir}
    lands in \code{BUILD\_RPATH} and \code{INSTALL\_RPATH}, so this usually
    means a dependency was not listed. Add it --- and remember the order is the
    link order.
  \item \textbf{A wasm target linking a native library.} Every \field{user\_lib}
    entry is taken as \code{lib<name>.a} built with the \emph{same} emsdk. A
    native \code{.so} cannot enter a wasm module.
  \item \textbf{Something the library needs was never initialised.} That is what
    \field{module\_init} is for. \S\ref{sec:moduleinit}
\end{itemize}

\section{The pipeline produced stale output}

\begin{itemize}
  \item \textbf{A generator binary from a previous emit.} If the emitted sources
    changed but the old \code{generator} still sits next to the output
    directory, running it re-emits the old bindings. The tool prints a rebuild
    note --- believe it.
  \item \textbf{A \field{user\_sources} glob.} Globs expand when
    \code{rosetta\_gen} runs. A new \code{.cpp} does not appear by itself.
  \item \textbf{Stale CMake caches} after a toolchain or layout change. Use
    \code{--fresh}.
\end{itemize}

\section{\code{--build} skipped a backend}

That is by design: a missing toolchain is a skip with a note, not an error.
The summary line says which tool was not found.

\begin{lstlisting}[style=plain]
   wasm        SKIPPED (emcc not found -- activate emsdk)
\end{lstlisting}

A backend that \emph{failed} is reported differently, and \code{--build} exits
non-zero after trying the rest.

\section{The wrong Python was used}

\code{--cmake-arg -DPython\_EXECUTABLE=\ldots} \textbf{does not work} --- the
generated CMake sets that variable with \code{CACHE \ldots\ FORCE} after
probing. Pin the interpreter on the target instead:

\begin{lstlisting}[style=json]
{ "lang": "nanobind", "python": "3.11", "requires_python": ">=3.10" }
\end{lstlisting}

This also fixes a wheel tagged for the wrong version, since the wheel is built
by whichever interpreter runs \code{make\_wheel.py}.

\section{\code{--clean} did not remove something}

By design in three cases: the manifest is never removed; a folder whose
\code{CMakeLists.txt} lacks the \code{Generated by rosetta} stamp is never
touched; and only the per-backend folders \emph{the manifest declares} are
removed, so anything else parked under \code{bindings/} survives.

\begin{lstlisting}[style=plain]
not removing gen (no rosetta_gen marker in its CMakeLists)
\end{lstlisting}

\section{A variable in the manifest did not expand}

Only a name declared in \field{variables} is substituted. \code{\$ENV\{HOME\}}
and any undeclared \code{\$\{\ldots\}} are CMake's and pass through untouched
--- and a \emph{misspelt} name is indistinguishable from those, so it passes
through silently and surfaces as a path that does not exist.

\code{\$\{NAME\}} rather than \code{\$NAME} is worth preferring: \code{\$Px}
is one name, \code{Px}.

\section{Reporting something}

The two artifacts that make a report actionable are \code{bindings/coverage.json}
and the manifest. Between them they say what rosetta thought your API was and
what it decided to do about it, which is almost always where the disagreement
is.
